\documentclass[11pt]{article}
\usepackage[margin=1in]{geometry}
\usepackage{amsmath,amssymb,amsthm}
\usepackage{hyperref}
\usepackage{enumitem}
\usepackage{xurl}
\let\oldus\_
\renewcommand{\_}{\oldus\allowbreak}

\newtheorem{theorem}{Theorem}
\newtheorem{lemma}[theorem]{Lemma}
\newtheorem{corollary}[theorem]{Corollary}
\theoremstyle{definition}
\newtheorem{definition}[theorem]{Definition}
\theoremstyle{remark}
\newtheorem{remark}[theorem]{Remark}

\title{A disproof of Conjecture 00000000223\\[2pt]
\large The density of $\{n : q \mid E_n\}$ is $0$, not a value of $q^{-2}$ type}
\date{}

\begin{document}
\sloppy
\maketitle

\section{The conjecture}

Conjecture 00000000223 reads:
\begin{quote}
\emph{Definition:} Euclid numbers $E_n = p_n\# + 1$ (one plus the product of the first $n$ primes).
\emph{Conjecture:} The least prime factor of $E_n$ exceeds $p_n$ uniformly for all sufficiently large $n$, and
the least prime factors are uniformly distributed on the logarithmic scale over $(p_n, p_n^2)$; moreover, for a
fixed prime $q$, the density of $n$ with $q$ dividing $E_n$ is an explicit value of $q^{-2}$ type.
\end{quote}
The Chinese text of the last clause says the density is a ``$q^{-2}$-type explicit value'' (a value of
$q^{-2}$ form). The conjecture is the conjunction of three clauses: (A) least prime factor $> p_n$; (B)
log-scale equidistribution of least prime factors; (C) the density clause. We refute (C), which refutes the
conjunction. Clauses (A) and (B) are not addressed. (Clause (A) is in fact true by the same argument, and (B)
is open; neither is used.)

\section{Definitions and conventions}

Following the conjecture and \url{https://en.wikipedia.org/wiki/Euclid_number} (``$E_n = p_n\# + 1$, where
$p_n\#$ is the nth primorial (the product of the first n prime numbers)''):
\begin{itemize}[nosep]
\item $p_1 = 2, p_2 = 3, \dots$; in Lean $p_{i+1} = \texttt{Nat.nth Nat.Prime } i$.
\item $p_n\# = \prod_{i<n} \texttt{nth Prime } i$ (Lean \texttt{primorialFirst}); Mathlib's \texttt{primorial}
  is the product of the primes $\le n$, a different function, and is not used.
\item $E_n = p_n\# + 1$ (Lean \texttt{euclid}); $D_q = \{n \in \mathbb{N} : q \mid E_n\}$ (Lean \texttt{divSet}).
\item $S \subseteq \mathbb{N}$ has natural density $d$ if $\#(S \cap [0,N))/N \to d$ (Lean
  \texttt{HasNatDensity}). We use $0 \in \mathbb{N}$; $E_0 = 2$. Adding or removing the single index $n=0$
  does not change any density.
\end{itemize}

\section{Readings of ``an explicit value of $q^{-2}$ type''}

We refute every reading that requires the density to be nonzero for at least one prime $q$, in particular:
\begin{enumerate}[label=(\roman*),nosep]
\item exact form: the density equals $c\,q^{-2}$ for every prime $q$, for some constant $c \neq 0$;
\item asymptotic form: the densities $d(q)$ exist and $q^2 d(q) \to c \neq 0$ as $q \to \infty$ through the
  primes (this is $d(q) \sim c\,q^{-2}$, and covers explicit values such as $1/(q^2-1)$).
\end{enumerate}
\textbf{Not refuted:} a pure upper-bound reading ``the density is $O(q^{-2})$''. The value $0$ satisfies it.
We do not regard it as the literal reading: the text asks for an explicit \emph{value} of $q^{-2}$ type
(form), not a bound, and $0$ has no dependence on $q$ at all. We state this reading explicitly so that the
reviewer can weigh it.

\section{The proof}

\begin{lemma}[Euclid; Lean \texttt{not\_dvd\_euclid}]
If $r < n$ then $p_{r+1} \nmid E_n$.
\end{lemma}
\begin{proof}
$p_{r+1}$ is one of the factors of $p_n\#$. If it also divided $p_n\# + 1$ it would divide $1$; but it is
prime. (Compare \url{https://en.wikipedia.org/wiki/Euclid%27s_theorem}: ``If p divides P and q, then p must
also divide the difference of the two numbers, which is (P + 1) - P or just 1.'')
\end{proof}

\begin{lemma}[Lean \texttt{divSet\_subset}, \texttt{hasNatDensity\_zero\_of\_subset}]
For a prime $q$ let $r = \#\{\text{primes} < q\}$, so $q = p_{r+1}$ (Mathlib \texttt{Nat.nth\_count}). Then
$D_q \subseteq [0, r]$, hence $\#(D_q \cap [0,N)) \le r+1$ and $D_q$ has natural density $0$.
\end{lemma}
\begin{proof}
The inclusion is the previous lemma. Then $0 \le \#(D_q \cap [0,N))/N \le (r+1)/N \to 0$.
\end{proof}

\begin{theorem}[Lean \texttt{conjecture\_223\_density\_clause\_false}]
\leavevmode
\begin{enumerate}[nosep]
\item For every prime $q$, $D_q$ has natural density $0$ (\texttt{hasNatDensity\_divSet}).
\item There is no prime $q$ and $d \ne 0$ such that $D_q$ has density $d$ (\texttt{no\_nonzero\_density}).
\item Reading (i) is false: there is no $c \ne 0$ with density $c/q^2$ for all primes $q$
  (\texttt{not\_exact\_form}; it already fails at $q = 2$).
\item Reading (ii) is false: there are no $d(\cdot)$ and $c \ne 0$ with $D_q$ of density $d(q)$ for all primes
  $q$ and $p_{r+1}^2\, d(p_{r+1}) \to c$ as $r \to \infty$ (\texttt{not\_asymptotic\_form}).
\end{enumerate}
\end{theorem}
\begin{proof}
(1) is the lemma. Limits in $\mathbb{R}$ are unique, so any density of $D_q$ equals $0$; this gives (2), and
(3) at $q = 2$. For (4), $p_{r+1}^2 d(p_{r+1}) = 0$ for all $r$, so the limit is $0 \ne c$.
\end{proof}

Example: $E_6 = 30031 = 59 \cdot 509$, so $59 \mid E_6$, but $59 = p_{17}$ divides no $E_n$ with $n \ge 17$.

\section{Lean formalization}

File \texttt{lean/Conjecture223/Basic.lean} (129 lines; only import \texttt{Mathlib}). Main
theorem \texttt{C223.conjecture\_223\_density\_clause\_false}, statement as in the Theorem above. Build:
\texttt{lake build} succeeds; \texttt{\#print axioms} for the capstone and for \texttt{no\_nonzero\_density},
\texttt{not\_exact\_form}, \texttt{not\_asymptotic\_form} reports only \texttt{propext},
\texttt{Classical.choice}, \texttt{Quot.sound}. No \texttt{sorry}, no \texttt{native\_decide}, no new axioms.

\end{document}
